%%%%%%%%%%%%
%
% $Autor: Wings $
% $Datum: 2019-03-05 08:03:15Z $
% $Pfad: Functions.tex $
% $Version: 4250 $
% !TeX spellcheck = en_GB/de_DE
% !TeX encoding = utf8
% !TeX root = manual 
% !TeX TXS-program:bibliography = txs:///biber
%
%%%%%%%%%%%%

\chapter{Functions}

\section{Core Localization Features}

\subsection{Visual Odometry Mode}
\begin{itemize}
    \item Estimates robot motion from consecutive camera frames
    \item Output: Relative pose transformation (translation and rotation)
    \item Good for short-term trajectory (< 20 meters)
\end{itemize}
\begin{figure}[H]
	\centering
	\begin{tikzpicture}[
		feat/.style={rectangle, draw=purple!60, fill=purple!5, thick, text width=4cm, text centered, rounded corners, minimum height=2.5cm, align=center},
		node distance=1cm
		]
		\node[feat] (f1) {\textbf{Visual Odometry}\\Estimates motion from frames\\Output: relative pose\\Range: $<$20m};
		\node[feat, right=of f1] (f2) {\textbf{Mapping Mode}\\Builds spatial map\\Output: 2D grid / 3D map\\Reference for loop closure};
		\node[feat, right=of f2] (f3) {\textbf{Localization Against Map}\\Matches frame to known map\\Output: global position\\Recovers from drift};
	\end{tikzpicture}
	\caption{Comparison of the three core localization features and their outputs}
	\label{fig:localization-features}
\end{figure}

\subsection{Mapping Mode}
\begin{itemize}
    \item Builds spatial map of environment from visual features
    \item Creates reference database for loop closure detection
    \item Output: 2D occupancy grid or 3D feature map
\end{itemize}

\subsection{Localization Against Map}
\begin{itemize}
    \item Matches current camera frame against known map
    \item Estimates global position even without encoder data
    \item Enables recovery from encoder drift
\end{itemize}

\section{Output Data}

\begin{itemize}
    \item Current position (x, y, theta) in world coordinates
    \item Confidence estimate for position accuracy
    \item Trajectory history (points visited)
    \item Loop closure events (when revisiting areas)
\end{itemize}


